%=====================================================================
% Jexpresso benchmark suite — description of the test cases listed in
% run_benchmark_list. Compile with pdflatex.
%
% Figure placeholders are provided for each case: replace the
% \figplaceholder{...} boxes with \includegraphics{...} once the
% result figures are available.
%=====================================================================
\documentclass[11pt,a4paper]{article}

\usepackage[margin=2.5cm]{geometry}
\usepackage{amsmath,amssymb,amsfonts}
\usepackage{bm}
\usepackage{graphicx}
\usepackage{booktabs}
\usepackage{siunitx}
\usepackage{hyperref}
\usepackage{xcolor}

% ---------------------------------------------------------------
% Placeholder box for result figures. Replace with
% \includegraphics[width=0.8\textwidth]{<figure file>} when ready.
% ---------------------------------------------------------------
\newcommand{\figplaceholder}[1]{%
  \fbox{\parbox[c][7cm][c]{0.85\textwidth}{\centering\itshape #1}}}

\newcommand{\dd}{\mathrm{d}}
\newcommand{\bU}{\bm{u}}

\title{The Jexpresso Benchmark Suite:\\
       Governing Equations, Setup, and Results of the Test Cases in
       \texttt{run\_benchmark\_list}}
\author{Jexpresso}
\date{\today}

\begin{document}
\maketitle

\begin{abstract}
This document describes the benchmark test cases executed by the script
\texttt{run\_benchmark\_list} of the Jexpresso spectral-element code.
For each case we report the governing equations, the initial and boundary
conditions, the source terms, and the numerical setup (grid, polynomial
order, time integrator, dissipation model). Placeholder figures are left
for the results of each benchmark. The suite comprises twelve serial cases
(compressible Euler, wave equation, elliptic/Helmholtz problems,
advection--diffusion, and shallow water) and three MPI-parallel cases
(rising thermal with tracers, a 3D rising thermal, and 2D flow over an
idealized city).
\end{abstract}

\tableofcontents
\clearpage

%=====================================================================
\section{Introduction}
%=====================================================================
Jexpresso is a research spectral-element-method (SEM) code for the
numerical solution of one-, two-, and three-dimensional partial
differential equations on quadrilateral and hexahedral grids, with
optional semi-infinite element extensions built from scaled Laguerre
functions. The benchmark suite collected in \texttt{run\_benchmark\_list}
exercises the main equation sets of the code:

\begin{itemize}
  \item \textbf{CompEuler}: compressible Euler equations in total-energy
        or potential-temperature form (plus a linear wave system used as
        a smoke test);
  \item \textbf{Elliptic} and \textbf{Helmholtz}: steady second-order
        problems solved with the assembled SEM Laplacian and an iterative
        Krylov solver;
  \item \textbf{AdvDiff}: linear advection--diffusion of passive scalars
        and a linearized wave-train problem;
  \item \textbf{ShallowWater}: the nonlinear shallow-water equations with
        bathymetry and wetting/drying.
\end{itemize}

The serial cases are run as
\verb|Jexpresso.run_case("<EquationSet>", "<case>")|, while the parallel
cases are launched under MPI, e.g.
\begin{verbatim}
mpiexec -n 2 julia --project=. src/Jexpresso.jl CompEuler thetaTracers
mpiexec -n 2 julia --project=. src/Jexpresso.jl CompEuler 3d
mpiexec -n 3 julia --project=. src/Jexpresso.jl CompEuler city2d
\end{verbatim}

All cases use Legendre--Gauss--Lobatto (LGL) nodal basis functions of
order $N$ ($\texttt{nop}$) within each element. Cases labelled
\emph{Laguerre} append a semi-infinite element layer discretized with
scaled Laguerre functions of order $N_{\rm Lag}$
($\texttt{nop\_laguerre}$) to absorb outgoing waves or to represent
unbounded directions.

%=====================================================================
\section{Compressible Euler equations (\texttt{CompEuler})}
%=====================================================================

\subsection{Governing equations}
Two formulations of the compressible Euler equations are used in the
suite. The \emph{total energy} form, used by the Sod shock-tube problem,
reads
\begin{equation}
  \frac{\partial}{\partial t}
  \begin{bmatrix} \rho \\ \rho u \\ \rho E \end{bmatrix}
  +
  \frac{\partial}{\partial x}
  \begin{bmatrix} \rho u \\ \rho u^2 + p \\ (\rho E + p)\,u \end{bmatrix}
  = \bm{0},
  \qquad
  p = (\gamma - 1)\,\rho\!\left(E - \tfrac{1}{2}u^2\right),
  \label{eq:euler-totalenergy}
\end{equation}
with $\gamma = c_p/c_v = 1.4$.

The atmospheric cases use the \emph{potential temperature}
(conservative $\rho\theta$) form. In two dimensions, with $\bm{q} =
(\rho,\, \rho u,\, \rho w,\, \rho\theta)^{T}$,
\begin{equation}
  \frac{\partial \bm{q}}{\partial t}
  + \frac{\partial \bm{F}}{\partial x}
  + \frac{\partial \bm{G}}{\partial z}
  = \bm{S} + \nabla\cdot(\mu \nabla \bm{q}),
  \label{eq:euler-theta}
\end{equation}
\begin{equation}
  \bm{F} =
  \begin{bmatrix} \rho u \\ \rho u^2 + p \\ \rho u w \\ \rho \theta u
  \end{bmatrix},
  \qquad
  \bm{G} =
  \begin{bmatrix} \rho w \\ \rho u w \\ \rho w^2 + p \\ \rho \theta w
  \end{bmatrix},
  \qquad
  \bm{S} =
  \begin{bmatrix} 0 \\ 0 \\ -\rho g \\ 0 \end{bmatrix},
\end{equation}
closed by the perfect-gas equation of state written in terms of
$\rho\theta$,
\begin{equation}
  p = p_{\rm ref}\left( \frac{\rho R \theta}{p_{\rm ref}}
      \right)^{c_p/c_v},
  \label{eq:eos-theta}
\end{equation}
where $R = c_p - c_v$ is the gas constant of dry air, $g =
\SI{9.81}{m.s^{-2}}$, and $p_{\rm ref} = \SI{e5}{Pa}$. The right-hand
side dissipation $\nabla\cdot(\mu\nabla\bm{q})$ stands for the selected
artificial viscosity / subgrid-scale (SGS) model (constant-coefficient
AV, dynamic SGS, Smagorinsky, or Vreman, depending on the case).
The hydrostatically balanced background used by the atmospheric cases is
\begin{equation}
  \theta_e = \theta_{\rm ref},
  \qquad
  p_e(z) = p_{\rm ref}\left( 1 - \frac{g z}{c_p\,\theta_{\rm ref}}
           \right)^{c_p/R},
  \qquad
  \rho_e = \rho(p_e, \theta_e),
  \label{eq:hydrostatic}
\end{equation}
and the perturbation variables $\bm{q}' = \bm{q} - \bm{q}_e$ are evolved
when the case is configured in perturbation (\texttt{PERT}) form.

%---------------------------------------------------------------------
\subsection{\texttt{sod1d}: Sod shock tube}
%---------------------------------------------------------------------
The classical Riemann problem of Sod for the 1D Euler equations in total
energy form~\eqref{eq:euler-totalenergy}. The initial condition is the
piecewise-constant state
\begin{equation}
  (\rho, u, p)(x, 0) =
  \begin{cases}
    (1.0,\ 0,\ 1.0)     & x < 0.5,\\[2pt]
    (0.125,\ 0,\ 0.1)   & x \ge 0.5,
  \end{cases}
  \qquad
  \rho E = \frac{p}{\gamma - 1} + \frac{1}{2}\rho u^2 .
\end{equation}
The exact solution consists of a right-moving shock, a contact
discontinuity, and a left-moving rarefaction wave.

\paragraph{Boundary conditions.} Dirichlet conditions holding the left
and right initial states at $x = 0$ and $x = 1$.

\paragraph{Numerical setup.}
\begin{center}
\begin{tabular}{ll}
\toprule
Domain                & $x \in [0, 1]$ \\
Elements              & $n_{\rm el} = 100$ \\
Polynomial order      & $N = 4$ (LGL) \\
Time integrator       & SSPRK53 \\
Time step             & $\Delta t = 10^{-4}$ \\
Final time            & $t_{\rm end} = 0.2$ \\
Stabilization         & Dynamic SGS (DSGS), $\mu = (1, 1, 1)$ \\
\bottomrule
\end{tabular}
\end{center}

\begin{figure}[htb!]
  \centering
  \figplaceholder{Placeholder: \texttt{sod1d} --- density, velocity, and
  pressure at $t = 0.2$ compared with the exact Riemann solution.}
  \caption{Sod shock tube: numerical solution at $t=0.2$.}
  \label{fig:sod1d}
\end{figure}

%---------------------------------------------------------------------
\subsection{\texttt{wave1d}: 1D linear wave system}
%---------------------------------------------------------------------
A first-order linear wave system for $(u, v)$,
\begin{equation}
  \frac{\partial u}{\partial t} + \frac{\partial v}{\partial x} = 0,
  \qquad
  \frac{\partial v}{\partial t} + \frac{\partial u}{\partial x} = 0,
  \label{eq:wave1d}
\end{equation}
equivalent to the second-order wave equation $u_{tt} = u_{xx}$ with unit
wave speed. The initial condition is a Gaussian-like pulse,
\begin{equation}
  u(x, 0) = 2^{-(x - 1)^2/\sigma^2},
  \qquad
  v(x, 0) = 0,
  \qquad \sigma = 0.15,
\end{equation}
which splits into left- and right-travelling pulses of half amplitude.

\paragraph{Boundary conditions.} Periodic in $x$.

\paragraph{Numerical setup.}
\begin{center}
\begin{tabular}{ll}
\toprule
Domain                & $x \in [0, 5]$, periodic \\
Elements              & $n_{\rm el} = 50$ \\
Polynomial order      & $N = 6$ (LGL) \\
Time integrator       & SSPRK33 \\
Time step             & $\Delta t = 10^{-3}$ \\
Final time            & $t_{\rm end} = 3.0$ \\
Dissipation           & none ($\texttt{lvisc} = \texttt{false}$) \\
\bottomrule
\end{tabular}
\end{center}

\begin{figure}[htb!]
  \centering
  \figplaceholder{Placeholder: \texttt{wave1d} --- $u(x,t)$ at selected
  times showing the counter-propagating pulses on the periodic domain.}
  \caption{1D wave system: evolution of the initial Gaussian pulse.}
  \label{fig:wave1d}
\end{figure}

%---------------------------------------------------------------------
\subsection{\texttt{wave1d\_lag}: 1D wave system with semi-infinite
Laguerre layers}
%---------------------------------------------------------------------
The same wave system~\eqref{eq:wave1d}, but on a finite domain
$x \in [-2.5, 2.5]$ extended on \emph{both} sides by semi-infinite
Laguerre elements that act as absorbing layers for outgoing waves. The
initial pulse is centered at the origin,
\begin{equation}
  u(x, 0) = 2^{-x^2/\sigma^2},
  \qquad
  v(x, 0) = 0,
  \qquad \sigma = 0.15 .
\end{equation}
In the absorbing layers, Rayleigh-type sponge terms are added,
\begin{equation}
  S_u = -c_s(x)\, u,
  \qquad
  S_v = -c_s(x)\, v,
  \qquad
  c_s(x) = \min\!\bigl(\beta(x),\, 1\bigr),
\end{equation}
where $\beta(x)$ is a smooth sigmoid ramp activated for $|x| \gtrsim
2.49$ (i.e., inside the Laguerre extensions only). The exact solution of
the truncated problem is a pair of pulses that leave the physical domain
without reflection; at long times the physical domain should return to
rest.

\paragraph{Boundary conditions.} Homogeneous Dirichlet ($u = 0$) at the
far ends of the Laguerre layers.

\paragraph{Numerical setup.}
\begin{center}
\begin{tabular}{ll}
\toprule
Physical domain       & $x \in [-2.5, 2.5]$ \\
Elements              & $n_{\rm el} = 50$ \\
Polynomial order      & $N = 6$ (LGL) \\
Laguerre extension    & both sides, $N_{\rm Lag} = 24$,
                        $\beta_{\rm Lag} = 1.0$ \\
Time integrator       & SSPRK33 \\
Time step             & $\Delta t = 10^{-3}$ \\
Final time            & $t_{\rm end} = 9.0$ \\
Dissipation           & none (sponge only) \\
\bottomrule
\end{tabular}
\end{center}

\begin{figure}[htb!]
  \centering
  \figplaceholder{Placeholder: \texttt{wave1d\_lag} --- $u(x,t)$ at
  selected times showing the pulses absorbed by the Laguerre layers
  without spurious reflection.}
  \caption{1D wave system with Laguerre absorbing layers.}
  \label{fig:wave1d-lag}
\end{figure}

%---------------------------------------------------------------------
\subsection{\texttt{theta}: 2D rising thermal bubble}
%---------------------------------------------------------------------
The classical warm-bubble benchmark for the potential-temperature Euler
equations~\eqref{eq:euler-theta}. A cone-shaped potential temperature
anomaly is superposed on the hydrostatically balanced, neutrally
stratified background~\eqref{eq:hydrostatic} with
$\theta_{\rm ref} = \SI{300}{K}$:
\begin{equation}
  \theta'(x, z, 0) =
  \begin{cases}
    \theta_c \left( 1 - \dfrac{r}{r_0} \right) & r < r_0,\\[6pt]
    0 & \text{otherwise},
  \end{cases}
  \qquad
  r = \sqrt{(x - x_c)^2 + (z - z_c)^2},
  \label{eq:bubble}
\end{equation}
with $\theta_c = \SI{2}{K}$, $r_0 = \SI{2000}{m}$, the bubble centered at
mid-domain horizontally and at $z_c = \SI{2500}{m}$, and $u = w = 0$
initially. Buoyancy lifts the bubble, which deforms into the
characteristic mushroom shape by $t \approx \SI{1000}{s}$.

\paragraph{Boundary conditions.} Free-slip (zero normal momentum) on all
boundaries:
$\;\rho\bU\cdot\bm{n} = 0$.

\paragraph{Numerical setup.}
\begin{center}
\begin{tabular}{ll}
\toprule
Grid                  & Gmsh TFI grid $10\times10$ elements
                        (\texttt{hexa\_TFI\_10x10.msh}) \\
Polynomial order      & $N = 4$ (LGL) \\
Form                  & perturbation (\texttt{PERT}) \\
Time integrator       & Carpenter--Kennedy LSRK(5,4) \\
Time step             & $\Delta t = \SI{0.5}{s}$ \\
Final time            & $t_{\rm end} = \SI{1000}{s}$ \\
Stabilization         & constant artificial viscosity,
                        $\mu = (0, 125, 125, 125)\,\si{m^2.s^{-1}}$ \\
\bottomrule
\end{tabular}
\end{center}

\begin{figure}[htb!]
  \centering
  \figplaceholder{Placeholder: \texttt{theta} --- potential temperature
  perturbation $\theta'$ at $t = \SI{1000}{s}$ (mushroom-shaped rising
  thermal).}
  \caption{2D rising thermal bubble: $\theta'$ at the final time.}
  \label{fig:theta}
\end{figure}

%---------------------------------------------------------------------
\subsection{\texttt{theta\_laguerre}: rising thermal bubble with
Laguerre top layer}
%---------------------------------------------------------------------
The same rising-bubble configuration as \texttt{theta}
(Eq.~\eqref{eq:bubble}, $\theta_c = \SI{2}{K}$, $r_0 = \SI{2000}{m}$,
$z_c = \SI{2500}{m}$, $\theta_{\rm ref} = \SI{300}{K}$), but the upper
boundary of the domain is replaced by a semi-infinite layer of scaled
Laguerre elements so that vertically propagating waves leave the domain
without reflection. In addition, a weak sponge is active above
$z_s = \SI{8000}{m}$:
\begin{equation}
  S_i \mathrel{-}= c_s\, q_i,
  \qquad
  c_s = \min\!\left( \sin^2\!\frac{\pi}{2}
        \frac{(z - z_s)}{z_{\rm max} - z_s},\, 1 \right),
  \qquad z \ge z_s .
\end{equation}

\paragraph{Boundary conditions.} Free-slip on the physical boundaries;
the Laguerre layer is attached at the boundary tagged
\texttt{free\_slip}.

\paragraph{Numerical setup.}
\begin{center}
\begin{tabular}{ll}
\toprule
Grid                  & Gmsh TFI grid $20\times20$
                        (\texttt{hexa\_TFI\_RTB20x20\_Lag.msh}) \\
Polynomial order      & $N = 4$ (LGL),
                        $N_{\rm Lag} = 16$ (vertical Laguerre) \\
Laguerre scaling      & $y_{\rm fac} = 110$ \\
Form                  & perturbation (\texttt{PERT}) \\
Time integrator       & Carpenter--Kennedy LSRK(5,4) \\
Time step             & $\Delta t = \SI{0.2}{s}$ \\
Final time            & $t_{\rm end} = \SI{1000}{s}$ \\
Stabilization         & $\mu = (0, 125, 125, 125)\,\si{m^2.s^{-1}}$;
                        erf-based filter, $\mu_x = \mu_y = 0.01$ \\
\bottomrule
\end{tabular}
\end{center}

\begin{figure}[htb!]
  \centering
  \figplaceholder{Placeholder: \texttt{theta\_laguerre} --- $\theta'$ at
  $t = \SI{1000}{s}$ with the semi-infinite Laguerre layer at the top of
  the domain.}
  \caption{Rising thermal bubble with Laguerre absorbing top layer.}
  \label{fig:theta-laguerre}
\end{figure}

%---------------------------------------------------------------------
\subsection{\texttt{thetaTracers}: rising thermal bubble with passive
tracers (MPI, 2 ranks)}
%---------------------------------------------------------------------
The potential-temperature system~\eqref{eq:euler-theta} augmented with
two passively advected tracers $q_{\rm tr}$ and $q_{\rm tr,2}$,
\begin{equation}
  \frac{\partial q_{\rm tr}}{\partial t}
  + \frac{\partial (q_{\rm tr} u)}{\partial x}
  + \frac{\partial (q_{\rm tr} w)}{\partial z} = 0,
\end{equation}
for a total of six unknowns
$(\rho, \rho u, \rho w, \rho\theta, q_{\rm tr}, q_{\rm tr,2})$.
The thermal anomaly is the cone~\eqref{eq:bubble} with
$\theta_c = \SI{2}{K}$, $r_0 = \SI{2000}{m}$, $z_c = \SI{2500}{m}$.
Tracer~1 is set to $1$ inside the bubble radius ($r < r_0$) and $0$
outside; tracer~2 is set to $1$ inside a second, smaller circle of
radius $r_{0,2} = \SI{1000}{m}$ centered at
$(x_{c,2}, z_{c,2}) = (-2000, 4000)\,\si{m}$. The case is run on an
\emph{unstructured} quadrilateral grid distributed over two MPI ranks.

\paragraph{Boundary conditions.} Free-slip on all boundaries.

\paragraph{Numerical setup.}
\begin{center}
\begin{tabular}{ll}
\toprule
Grid                  & unstructured Gmsh grid
                        (\texttt{square\_UNSTR\_20el.msh}), 2 MPI ranks \\
Polynomial order      & $N = 4$ (LGL) \\
Time integrator       & Carpenter--Kennedy LSRK(5,4) \\
Time step             & $\Delta t = \SI{0.2}{s}$ \\
Final time            & $t_{\rm end} = \SI{1000}{s}$ \\
Stabilization         & Smagorinsky SGS,
                        $\mu = (0, 1, 1, 2, 3, 1)$ \\
\bottomrule
\end{tabular}
\end{center}

\begin{figure}[htb!]
  \centering
  \figplaceholder{Placeholder: \texttt{thetaTracers} --- $\theta'$ and
  tracer concentrations $q_{\rm tr}$, $q_{\rm tr,2}$ at
  $t = \SI{1000}{s}$ on the unstructured grid (2 MPI ranks).}
  \caption{Rising thermal with two passive tracers.}
  \label{fig:thetaTracers}
\end{figure}

%---------------------------------------------------------------------
\subsection{\texttt{3d}: 3D rising thermal bubble (MPI, 2 ranks)}
%---------------------------------------------------------------------
The three-dimensional extension of the rising-bubble problem, with
unknowns $(\rho, \rho u, \rho v, \rho w, \rho\theta)$ and fluxes
\begin{equation}
  \bm{F} =
  \begin{bmatrix} \rho u \\ \rho u^2 + p \\ \rho u v \\ \rho u w \\
                  \rho\theta u \end{bmatrix},
  \quad
  \bm{G} =
  \begin{bmatrix} \rho v \\ \rho u v \\ \rho v^2 + p \\ \rho v w \\
                  \rho\theta v \end{bmatrix},
  \quad
  \bm{H} =
  \begin{bmatrix} \rho w \\ \rho u w \\ \rho v w \\ \rho w^2 + p \\
                  \rho\theta w \end{bmatrix},
  \quad
  \bm{S} =
  \begin{bmatrix} 0 \\ 0 \\ 0 \\ -\rho g \\ 0 \end{bmatrix}.
\end{equation}
The thermal anomaly is again Eq.~\eqref{eq:bubble} with
$\theta_c = \SI{2}{K}$, $r_0 = \SI{2000}{m}$, applied in the $x$--$z$
plane with bubble center height $z_c = \SI{2500}{m}$ (the domain is one
element thick in $y$, making this a quasi-2D verification of the 3D MPI
infrastructure).

\paragraph{Boundary conditions.} Free-slip on all boundaries.

\paragraph{Numerical setup.}
\begin{center}
\begin{tabular}{ll}
\toprule
Grid                  & Gmsh TFI grid $10\times1\times10$
                        (\texttt{hexa\_TFI\_10x1x10.msh}), 2 MPI ranks \\
Polynomial order      & $N = 4$ (LGL) \\
Time integrator       & Carpenter--Kennedy LSRK(5,4) \\
Time step             & $\Delta t = \SI{0.5}{s}$ \\
Final time            & $t_{\rm end} = \SI{1000}{s}$ \\
Stabilization         & Vreman SGS model, $\mu = (0, 1, 1, 1, 2)$ \\
\bottomrule
\end{tabular}
\end{center}

\begin{figure}[htb!]
  \centering
  \figplaceholder{Placeholder: \texttt{3d} --- iso-surfaces / mid-plane
  slice of $\theta'$ at $t = \SI{1000}{s}$ from the 3D MPI run.}
  \caption{3D rising thermal bubble.}
  \label{fig:3d}
\end{figure}

%---------------------------------------------------------------------
\subsection{\texttt{city2d}: 2D flow over an idealized city (MPI, 3
ranks)}
%---------------------------------------------------------------------
Neutral atmospheric flow past a set of rectangular building obstacles in
a $\SI{1000}{m} \times \SI{1000}{m}$ vertical slice, governed by the
potential-temperature Euler equations~\eqref{eq:euler-theta} in
\emph{total} form. The atmosphere is initialized as the neutrally
stratified hydrostatic state~\eqref{eq:hydrostatic} with
$\theta_{\rm ref} = \SI{300}{K}$ and a uniform horizontal wind
\begin{equation}
  u(x, z, 0) = \SI{10}{m.s^{-1}},
  \qquad
  w(x, z, 0) = 0 .
\end{equation}
A Rayleigh sponge relaxing the solution toward the background state
$\bm{q}_e$ is active above $z_s = \SI{500}{m}$,
\begin{equation}
  S_i \mathrel{-}= c_s\,(q_i - q_{e,i}),
  \qquad
  c_s = \alpha \sin\!\frac{\pi}{2}\frac{z - z_s}{z_{\rm max} - z_s},
  \qquad \alpha = 0.05,
  \quad z \ge z_s,
\end{equation}
for the momentum and $\rho\theta$ equations. The flow separates at the
building rooftops and produces unsteady recirculation in the street
canyons.

\paragraph{Boundary conditions.} Free-slip on the ground, on the
building walls, and on the outer boundaries.

\paragraph{Numerical setup.}
\begin{center}
\begin{tabular}{ll}
\toprule
Grid                  & Gmsh transfinite grid with embedded buildings
                        (\texttt{city2d\_transfinite.msh}), 3 MPI ranks \\
Polynomial order      & $N = 4$ (LGL) \\
Form                  & total (\texttt{TOTAL}) \\
Time integrator       & Carpenter--Kennedy LSRK(5,4) \\
Time step             & $\Delta t = \SI{0.004}{s}$ \\
Final time            & $t_{\rm end} = \SI{600}{s}$ \\
Stabilization         & Smagorinsky SGS, $\mu = (0, 5, 5, 5)$ \\
Sponge                & $z_s = \SI{500}{m}$, $\alpha = 0.05$ \\
\bottomrule
\end{tabular}
\end{center}

\begin{figure}[htb!]
  \centering
  \figplaceholder{Placeholder: \texttt{city2d} --- horizontal velocity
  $u$ (and/or vorticity) around the buildings at $t = \SI{600}{s}$ from
  the 3-rank MPI run.}
  \caption{2D flow over an idealized city.}
  \label{fig:city2d}
\end{figure}

\clearpage
%=====================================================================
\section{Elliptic problem (\texttt{Elliptic/case1})}
%=====================================================================
A steady Poisson/Laplace problem for a scalar $u(x, y)$,
\begin{equation}
  \nabla^2 u = f \quad \text{in } \Omega,
  \qquad f = 0,
\end{equation}
on the rectangle $\Omega = [1, 6] \times [0, 3.14]$ (a unit square mesh
scaled by $(L_x, L_y) = (5.0,\, 3.14)$ and displaced by $1$ in $x$).
With $f = 0$ this is the Laplace equation driven entirely by the
Dirichlet data:
\begin{equation}
  u = 100 \ \text{on the left, right, and bottom boundaries},
  \qquad
  u = 100 + 20 \sin\!\frac{\pi x}{L_x} \ \text{on the top boundary}.
\end{equation}
The SEM stiffness (Laplacian) matrix is assembled in weak form,
$\int_\Omega \nabla\phi_i \cdot \nabla\phi_j \,\dd\Omega$, with
continuity imposed by direct stiffness summation (DSS), and the linear
system is solved with the BiCGStab Krylov method.

\paragraph{Numerical setup.}
\begin{center}
\begin{tabular}{ll}
\toprule
Grid                  & Gmsh grid \texttt{square\_dirichletT.msh},
                        scaled to $[1,6]\times[0,3.14]$ \\
Polynomial order      & $N = 4$ (LGL) \\
Linear solver         & BiCGStab (\texttt{BICGSTABLE}) \\
Initial guess         & $u = 0$ \\
\bottomrule
\end{tabular}
\end{center}

\begin{figure}[htb!]
  \centering
  \figplaceholder{Placeholder: \texttt{Elliptic/case1} --- steady
  solution $u(x,y)$ of the Laplace problem with the sinusoidally heated
  top boundary.}
  \caption{Elliptic (Laplace) problem: converged solution.}
  \label{fig:elliptic-case1}
\end{figure}

%=====================================================================
\section{Helmholtz problems (\texttt{Helmholtz})}
%=====================================================================

\subsection{\texttt{case1}: Helmholtz equation with manufactured
solution}
A Helmholtz problem with reaction coefficient $\alpha = 10$,
\begin{equation}
  \nabla^2 u + \alpha\, u = f \quad \text{in } \Omega
  = [1, 6] \times [0, 3.14],
  \qquad u = 0 \ \text{on } \partial\Omega,
  \label{eq:helmholtz}
\end{equation}
where the right-hand side is manufactured from the exact solution
\begin{equation}
  u_e(x, y) = -\sin\!\frac{x}{2}\; e^{-x/2} \cos y,
\end{equation}
so that the discrete source assembled by the code is
\begin{equation}
  S(x, y) = \frac{1}{2}\cos\!\frac{x}{2}\, e^{-x/2}\cos y
            + (1 + \alpha) \sin\!\frac{x}{2}\, e^{-x/2}\cos y .
\end{equation}
The error against $u_e$ measures the spectral accuracy of the assembled
SEM Helmholtz operator.

\paragraph{Numerical setup.}
\begin{center}
\begin{tabular}{ll}
\toprule
Grid                  & Gmsh TFI grid \texttt{hexa\_TFI\_helmholtz.msh},
                        scaled to $[1,6]\times[0,3.14]$ \\
Polynomial order      & $N = 10$ (LGL) \\
Helmholtz coefficient & $\alpha = 10$ \\
Linear solver         & BiCGStab (\texttt{BICGSTABLE}) \\
Boundary conditions   & homogeneous Dirichlet \\
\bottomrule
\end{tabular}
\end{center}

\begin{figure}[htb!]
  \centering
  \figplaceholder{Placeholder: \texttt{Helmholtz/case1} --- computed
  $u(x,y)$ and pointwise error against the exact solution
  $u_e = -\sin(x/2)\,e^{-x/2}\cos y$.}
  \caption{Helmholtz problem on the bounded domain.}
  \label{fig:helmholtz-case1}
\end{figure}

\subsection{\texttt{case1\_laguerre}: Helmholtz equation on a
semi-infinite domain}
The same manufactured Helmholtz problem~\eqref{eq:helmholtz}
($\alpha = 10$, same $u_e$, $f$, and homogeneous Dirichlet data), but
the domain is extended to a semi-infinite strip in the $x$-direction by
a layer of scaled Laguerre functions appended to the LGL grid. Since
$u_e \propto e^{-x/2}$ decays exponentially, the Laguerre basis
(which embeds exponential decay) is the natural discretization of the
unbounded direction; the case verifies the coupled LGL--Laguerre
assembly of the Helmholtz operator.

\paragraph{Numerical setup.}
\begin{center}
\begin{tabular}{ll}
\toprule
Grid                  & Gmsh grid
                        \texttt{hexa\_TFI\_helmholtz\_laguerre.msh}
                        + Laguerre extension in $x$ \\
Polynomial order      & $N = 10$ (LGL), $N_{\rm Lag} = 14$ (Laguerre) \\
Laguerre scaling      & $x_{\rm fac} = 0.25$, $\beta_{\rm Lag} = 1.0$ \\
Helmholtz coefficient & $\alpha = 10$ \\
Linear solver         & BiCGStab (\texttt{BICGSTABLE}) \\
\bottomrule
\end{tabular}
\end{center}

\begin{figure}[htb!]
  \centering
  \figplaceholder{Placeholder: \texttt{Helmholtz/case1\_laguerre} ---
  computed $u(x,y)$ on the LGL grid + Laguerre semi-infinite extension,
  and error against $u_e$.}
  \caption{Helmholtz problem with a semi-infinite Laguerre extension.}
  \label{fig:helmholtz-laguerre}
\end{figure}

\clearpage
%=====================================================================
\section{Advection--diffusion problems (\texttt{AdvDiff})}
%=====================================================================

\subsection{\texttt{Wave\_Train}: 1D linearized wave train with Laguerre
outflow}
A 1D linearized shallow-water (wave) system for the free-surface
perturbation $h$ and velocity $u$ over a still background of depth
$H = \SI{10}{m}$,
\begin{equation}
  \frac{\partial h}{\partial t}
  + \frac{\partial (H u)}{\partial x} = 0,
  \qquad
  \frac{\partial u}{\partial t}
  + \frac{\partial (g h)}{\partial x} = 0,
  \qquad g = \SI{9.81}{m.s^{-2}},
\end{equation}
(the general fluxes are $F_h = U h + H u$ and $F_u = g h + U u$ with
background current $U = 0$ here). The fluid starts at rest,
$h = u = 0$, and a monochromatic wave train is injected by the
time-dependent Dirichlet condition at the left boundary,
\begin{equation}
  u(0, t) = 0.025 \sin\!\left( \frac{2\pi\, 30\, t}{5000} \right)
  \,\si{m.s^{-1}},
\end{equation}
i.e., amplitude $\SI{0.025}{m.s^{-1}}$ and period
$T = 5000/30 \approx \SI{166.7}{s}$. The right end of the domain is
extended with a semi-infinite layer of Laguerre elements so that the
wave train exits without reflection.

\paragraph{Numerical setup.}
\begin{center}
\begin{tabular}{ll}
\toprule
Domain                & $x \in [0, 5000]\,\si{m}$ \\
Elements              & $n_{\rm el} = 300$ \\
Polynomial order      & $N = 4$ (LGL) \\
Laguerre extension    & right side, $N_{\rm Lag} = 20$,
                        $\beta_{\rm Lag} = 1.0$, $y_{\rm fac} = 100$ \\
Time integrator       & SSPRK33 \\
Time step             & $\Delta t = \SI{0.02}{s}$ \\
Final time            & $t_{\rm end} = \SI{5000}{s}$ \\
Dissipation           & none \\
\bottomrule
\end{tabular}
\end{center}

\begin{figure}[htb!]
  \centering
  \figplaceholder{Placeholder: \texttt{Wave\_Train} --- free-surface
  perturbation $h(x)$ at selected times showing the wave train
  propagating into the Laguerre outflow layer without reflection.}
  \caption{1D wave train with Laguerre radiation layer.}
  \label{fig:wavetrain}
\end{figure}

\subsection{\texttt{kopriva}: 2D advection--diffusion of a Gaussian
(Kopriva benchmark)}
The linear advection--diffusion equation for a scalar $q(x, y, t)$,
\begin{equation}
  \frac{\partial q}{\partial t}
  + \nabla\cdot(\bU q)
  = \nu \nabla^2 q,
  \qquad
  \bU = (u, v) = (0.5,\ 1.0),
  \qquad \nu = 0.1,
  \label{eq:advdiff}
\end{equation}
following the classical test of Kopriva for spectral element methods,
for which the exact solution is an advecting, spreading Gaussian. The
initial condition is
\begin{equation}
  q(x, y, 0) = \exp\!\left[ -\frac{(x - x_c)^2}{\sigma_x^2}
                            -\frac{(y - y_c)^2}{\sigma_y^2} \right],
  \qquad \sigma_x = \sigma_y = 1,
\end{equation}
centered at mid-domain in $x$ and at $y_c = 3$, on a periodic Gmsh grid.

\paragraph{Numerical setup.}
\begin{center}
\begin{tabular}{ll}
\toprule
Grid                  & periodic Gmsh grid
                        (\texttt{kopriva\_periodic.msh}) \\
Polynomial order      & $N = 4$ (LGL) \\
Time integrator       & SSPRK54 \\
Time step             & $\Delta t = \SI{0.005}{s}$ \\
Final time            & $t_{\rm end} = \SI{10}{s}$ \\
Diffusivity           & $\nu_x = \nu_y = 0.1$ \\
\bottomrule
\end{tabular}
\end{center}

\begin{figure}[htb!]
  \centering
  \figplaceholder{Placeholder: \texttt{kopriva} --- contours of $q$ at
  $t = 0.5, 1, 2, 4$ compared with the exact advecting--diffusing
  Gaussian.}
  \caption{Kopriva advection--diffusion benchmark.}
  \label{fig:kopriva}
\end{figure}

\subsection{\texttt{2d\_Laguerre}: 2D advection--diffusion with Laguerre
absorbing layers}
The same advection--diffusion equation~\eqref{eq:advdiff} with
$\bU = (0.5, 1.0)$ and $\nu = 0.1$, but on a bounded grid augmented with
Laguerre semi-infinite layers (in $y$) acting as an absorbing outflow
for the scalar. The initial condition is the Gaussian
\begin{equation}
  q(x, y, 0) = \exp\!\left[ -\frac{(x - x_c)^2}{\sigma_x^2}
                            -\frac{(y - y_c)^2}{\sigma_y^2} \right],
  \qquad \sigma_x = \sigma_y = 1,
\end{equation}
centered at mid-domain in $x$ and $y_c = 8$. Homogeneous Dirichlet
values are imposed on the physical inflow boundaries.

\paragraph{Numerical setup.}
\begin{center}
\begin{tabular}{ll}
\toprule
Grid                  & Gmsh grid \texttt{Wave\_Train.msh}
                        (scaled by $10\times10$) + Laguerre layer \\
Polynomial order      & $N = 4$ (LGL), $N_{\rm Lag} = 30$ \\
Laguerre scaling      & $y_{\rm fac} = 0.07$ \\
Time integrator       & SSPRK54 \\
Time step             & $\Delta t = \SI{5e-4}{s}$ \\
Final time            & $t_{\rm end} = \SI{4}{s}$ \\
Diffusivity           & $\nu = 0.1$ \\
\bottomrule
\end{tabular}
\end{center}

\begin{figure}[htb!]
  \centering
  \figplaceholder{Placeholder: \texttt{2d\_Laguerre} --- contours of $q$
  at $t = 0.5, 1, 2, 4$ showing the Gaussian advected into the Laguerre
  layer.}
  \caption{2D advection--diffusion with Laguerre absorbing layer.}
  \label{fig:2dlaguerre}
\end{figure}

\clearpage
%=====================================================================
\section{Shallow-water equations
(\texttt{ShallowWater/SoliWaveIsland})}
%=====================================================================
The nonlinear shallow-water equations in conservative form with
bathymetry source terms,
\begin{align}
  \frac{\partial H}{\partial t}
  + \frac{\partial (Hu)}{\partial x}
  + \frac{\partial (Hv)}{\partial y} &= 0, \\
  \frac{\partial (Hu)}{\partial t}
  + \frac{\partial}{\partial x}\!\left( Hu^2 + \tfrac{1}{2} g H^2 \right)
  + \frac{\partial (Huv)}{\partial y}
  &= -g H \frac{\partial H_b}{\partial x}, \\
  \frac{\partial (Hv)}{\partial t}
  + \frac{\partial (Huv)}{\partial x}
  + \frac{\partial}{\partial y}\!\left( Hv^2 + \tfrac{1}{2} g H^2 \right)
  &= -g H \frac{\partial H_b}{\partial y},
\end{align}
where $H$ is the water depth, $H_b(x,y)$ the bathymetry (bottom
elevation), and $g = \SI{9.81}{m.s^{-2}}$. The implementation uses a
well-balanced perturbation splitting about the lake-at-rest state
$H_e = \max(H_{\rm dry},\ h_0 - H_b)$: the pressure flux carries
$\tfrac{1}{2} g (H^2 - H_e^2)$ and the bathymetry source carries
$-g (H - H_e)\nabla H_b$, so that the lake at rest is preserved exactly.
Velocities are recovered from the momenta with the desingularization of
Kurganov \& Petrova,
\begin{equation}
  u = \frac{\sqrt{2}\, H\, (Hu)}{\sqrt{H^4 + \max(H, H_{\rm wet})^4}},
  \qquad H_{\rm wet} = 10^{-3}\,\si{m},
\end{equation}
and momentum in dry cells ($H < H_{\rm wet}$) is relaxed to zero at rate
$\sigma_{\rm dry} = \SI{25}{s^{-1}}$.

\paragraph{Setup: solitary wave impinging on a conical island.}
Following the laboratory configuration of a solitary wave interacting
with a conical island (cf.\ Synolakis-type solitary-wave initialization),
the bathymetry is the cone
\begin{equation}
  H_b(x, y) = h_c \max\!\left( 0,\ 1 - \frac{r}{r_c} \right),
  \qquad
  r = \sqrt{(x - x_{c})^2 + y^2},
  \qquad
  h_c = \SI{0.93}{m}, \quad r_c = \SI{3.6}{m},
\end{equation}
centered at $(x_c, y_c) = (12.5, 0)\,\si{m}$ in a basin of still-water
depth $h_0 = \SI{0.32}{m}$ (the cone pierces the free surface near its
center). The initial condition is the solitary wave
\begin{equation}
  \eta(x) = A\,\mathrm{sech}^2\!\bigl[\gamma (x - x_w)\bigr],
  \qquad
  \gamma = \sqrt{\frac{3A}{4 h_0^3}},
  \qquad
  A = \SI{0.064}{m}, \quad x_w = \SI{2.5}{m},
\end{equation}
with the associated velocity field
\begin{equation}
  u(x) = \frac{c\,\eta}{h_0 + \eta},
  \qquad
  c = \sqrt{g (h_0 + \eta)},
  \qquad v = 0,
\end{equation}
and depth $H(x, y, 0) = \max\bigl(H_{\rm dry},\ h_0 + \eta(x) -
H_b(x,y)\bigr)$, $H_{\rm dry} = 10^{-3}\,\si{m}$. The wave shoals on the
island flank, partially runs up, wraps around the island, and the two
wave fronts collide on the lee side.

\paragraph{Boundary conditions.} Free-slip reflective walls on all four
boundaries (normal momentum removed:
$(Hu, Hv) \leftarrow (Hu, Hv) - [(Hu, Hv)\cdot\bm{n}]\,\bm{n}$).

\paragraph{Numerical setup.}
\begin{center}
\begin{tabular}{ll}
\toprule
Domain                & $[0, 25] \times [-15, 15]\,\si{m}$ \\
Grid                  & Gmsh transfinite grid, $25 \times 30 = 750$
                        quadrilaterals (\texttt{SoliWaveIsland.msh}) \\
Polynomial order      & $N = 4$ (LGL) \\
Time integrator       & SSPRK54 \\
Time step             & $\Delta t = \SI{0.01}{s}$ \\
Final time            & $t_{\rm end} = \SI{25}{s}$ \\
Stabilization         & constant artificial viscosity,
                        $\mu = (0.05, 0.05, 0.05)\,\si{m^2.s^{-1}}$,
                        applied to $(H - H_e, Hu, Hv)$ \\
Wet/dry treatment     & $H_{\rm wet} = 10^{-3}\,\si{m}$,
                        $\sigma_{\rm dry} = \SI{25}{s^{-1}}$ \\
\bottomrule
\end{tabular}
\end{center}

\begin{figure}[htb!]
  \centering
  \figplaceholder{Placeholder: \texttt{SoliWaveIsland} --- free-surface
  elevation snapshots at selected times ($t \approx 4, 8, 12, 16, 20$
  \si{s}) showing the solitary wave running up and wrapping around the
  conical island.}
  \caption{Solitary wave interacting with a conical island.}
  \label{fig:soliwave}
\end{figure}

\clearpage
%=====================================================================
\section{Summary}
%=====================================================================
\begin{table}[htb!]
\centering
\small
\caption{Overview of the benchmark cases in
\texttt{run\_benchmark\_list}.}
\label{tab:summary}
\begin{tabular}{llllll}
\toprule
Case & Equations & Dim. & $N$ & Integrator & $t_{\rm end}$ \\
\midrule
\texttt{CompEuler/sod1d}          & Euler ($\rho E$)        & 1D & 4  & SSPRK53 & 0.2 \\
\texttt{CompEuler/wave1d}         & linear wave             & 1D & 6  & SSPRK33 & 3 \\
\texttt{CompEuler/wave1d\_lag}    & linear wave + Laguerre  & 1D & 6  & SSPRK33 & 9 \\
\texttt{CompEuler/theta}          & Euler ($\rho\theta$)    & 2D & 4  & LSRK(5,4) & 1000 \\
\texttt{CompEuler/theta\_laguerre}& Euler ($\rho\theta$) + Laguerre & 2D & 4 & LSRK(5,4) & 1000 \\
\texttt{Elliptic/case1}           & $\nabla^2 u = 0$        & 2D & 4  & BiCGStab & --- \\
\texttt{Helmholtz/case1}          & $\nabla^2 u + 10u = f$  & 2D & 10 & BiCGStab & --- \\
\texttt{Helmholtz/case1\_laguerre}& Helmholtz + Laguerre    & 2D & 10/14 & BiCGStab & --- \\
\texttt{AdvDiff/Wave\_Train}      & linear wave system      & 1D & 4  & SSPRK33 & 5000 \\
\texttt{AdvDiff/kopriva}          & advection--diffusion    & 2D & 4  & SSPRK54 & 10 \\
\texttt{AdvDiff/2d\_Laguerre}     & adv--diff + Laguerre    & 2D & 4/30 & SSPRK54 & 4 \\
\texttt{ShallowWater/SoliWaveIsland} & shallow water        & 2D & 4  & SSPRK54 & 25 \\
\midrule
\texttt{CompEuler/thetaTracers} (2 ranks) & Euler ($\rho\theta$) + 2 tracers & 2D & 4 & LSRK(5,4) & 1000 \\
\texttt{CompEuler/3d} (2 ranks)           & Euler ($\rho\theta$)             & 3D & 4 & LSRK(5,4) & 1000 \\
\texttt{CompEuler/city2d} (3 ranks)       & Euler ($\rho\theta$)             & 2D & 4 & LSRK(5,4) & 600 \\
\bottomrule
\end{tabular}
\end{table}

\end{document}
